% ---------------------------------------------------------------------------
% Methods: how the denoising error is measured, and why. Moved here from
% results.tex on 2026-10-14 and extended with the reasons for each choice.
% Numbers checked on 2026-10-14: 0.67% and 40.7 deg from 200,000 random pairs
% of cubic orientations (src/orientation.py); 44% from the uniform angle
% (P(angle > 5) = 4/9 at theta_max = 9 deg).
% ---------------------------------------------------------------------------

\subsection{Error measures}
\label{sec:metrics}

\paragraph{Disorientation.}
To compare a denoised map with the clean map we need the angle between two
orientations. For a cubic crystal the same orientation can be written in 24
equivalent ways, so a plain difference between two sets of Euler angles is not
a valid measure~\cite{atindama2023restoration,ding2020indexing}. We therefore
use the disorientation angle. Let $\mathbf{q}$ and $\mathbf{q}'$ be the unit
quaternions of two orientations. Their misorientation is
\begin{equation}
  \Delta\mathbf{q} = \mathbf{q}'\,\mathbf{q}^{-1},
  \label{eq:misorientation}
\end{equation}
and the disorientation angle is the smallest rotation angle among the
symmetrically equivalent forms of $\Delta\mathbf{q}$,
\begin{equation}
  \delta(\mathbf{q}, \mathbf{q}') = \min_{\mathbf{s} \in \mathcal{S}}
  \; 2 \arccos \left| \left( \mathbf{s}\, \Delta\mathbf{q} \right)_0 \right|,
  \label{eq:disorientation}
\end{equation}
where $\mathcal{S}$ is the set of the 24 proper rotations of the cubic point
group and $(\cdot)_0$ is the scalar part of a quaternion. Every error in this
paper is a disorientation angle.

\paragraph{Four measures.}
A single number cannot show everything a denoising method does to a map, so
we use four measures, each answering one question.

\emph{ALL: how large is the error overall?} For a map $m$ with the set of
pixels $\Omega$ ($|\Omega| = 128 \times 128 = 16{,}384$), denoised map $u_m$
and true map $u^0_m$, the error is the mean per-pixel disorientation,
\begin{equation}
  e_m = \frac{1}{|\Omega|} \sum_{\mathbf{x} \in \Omega}
        \delta\big(u_m(\mathbf{x}), u^0_m(\mathbf{x})\big).
  \label{eq:error}
\end{equation}
We use the mean, as Atindama et al.~\cite{atindama2023restoration} did, and
not the median. A median ignores any group of pixels smaller than half the
map, so it would not change at all if the misindexed pixels, at most 9\% of a
test map, were left unrepaired.

\emph{BOUNDARY: are grain boundaries kept sharp?} The same mean over the
grain-boundary pixels $B_m$, the pixels with at least one of their four
neighbours in a different grain of the true map. Since every DREAM.3D grain
has a single orientation, these pixels are found exactly, without a threshold
angle. About 17\% of the pixels lie on a boundary. This is where a filter
averages two grains and blurs the boundary, and where most misindexing
occurs. A method can lower ALL by smoothing the grain interiors, which hold
83\% of the pixels, while blurring the boundaries; BOUNDARY shows this.

\emph{CORRUPTED: were the misindexed pixels repaired?} The same mean over the
misindexed pixels $M_m$, which the noise generator replaced and recorded for
every map. A replaced pixel starts on average $40.7^\circ$ from its true
orientation, so this measure shows directly whether a method found and
rebuilt it. It is defined only for test sets with misindexing, and only for
synthetic data, where the replaced pixels are known.

\emph{${>}10^\circ$: how many pixels are still grossly wrong?} The percentage
of pixels whose error exceeds $10^\circ$,
\begin{equation}
  g_m = \frac{100}{|\Omega|} \sum_{\mathbf{x} \in \Omega}
        \mathbb{1}\Big[\delta\big(u_m(\mathbf{x}), u^0_m(\mathbf{x})\big) > 10^\circ\Big],
  \label{eq:gross}
\end{equation}
where $\mathbb{1}[\cdot]$ is 1 when the condition holds and 0 otherwise. A few
wrong pixels barely move the mean: if 1\% of the pixels are $60^\circ$ off and
the rest $0.3^\circ$, ALL is only $0.9^\circ$, yet about 160 pixels of the map
are wrong. We chose $10^\circ$ because it is a common threshold for a grain
boundary in EBSD analysis, and the smallest grain threshold we give the MTEX
filters (Section~\ref{sec:baselines}): a pixel more than $10^\circ$ off would be
taken for a separate grain, and so creates false grains and boundaries. A
smaller threshold, such as $5^\circ$, would count ordinary scatter as gross
error: at $\theta_{\max} = 9^\circ$, scatter alone moves 44\% of the pixels
beyond $5^\circ$ but none beyond $10^\circ$. Conversely, only 0.67\% of
random orientation pairs lie within $10^\circ$ of each other, so practically
every misindexed pixel that is not repaired is counted. At
$\theta_{\max} \geq 12^\circ$, scatter alone moves 17--33\% of the pixels
beyond $10^\circ$ (Section~\ref{sec:test_conditions}). At these levels the
measure mixes the two kinds of error, and CORRUPTED is the measure of
misindexing repair.

\paragraph{Averaging over maps.}
Each test set has $N = 500$ maps. We compute every measure per map first and
then report the mean over the maps, for example
$\bar{e} = \frac{1}{N}\sum_{m=1}^{N} e_m$, with the standard deviation across
the maps. The map, not the pixel, is the unit of every statistic: neighbouring
pixels share a grain and a filter window, so their errors are not independent,
and treating the 16{,}384 pixels of a map as independent samples would
overstate the certainty of a result. Since every map has the same number of
pixels, $\bar{e}$ is also the mean disorientation over all $N|\Omega|$ pixels
of the set.

\paragraph{Repeated training.}
We trained the network three times, with different random seeds. Unless
stated otherwise, the value given for our method is the mean of the three
runs, and the $\pm$ value is the standard deviation across the runs, which
shows how much a result depends on the random start of training. A
difference between two methods smaller than this spread is not a result.

\paragraph{Fraction of noise removed.}
The input error grows with $\theta_{\max}$, so absolute errors cannot be
compared between noise levels. We therefore also give the fraction of the
input error that the denoising removed,
\begin{equation}
  R = 1 - \frac{\bar{e}_{\text{out}}}{\bar{e}_{\text{in}}},
  \label{eq:removed}
\end{equation}
which puts all levels on the same scale.

\paragraph{Statistical test.}
To test whether two methods differ, we use the paired $t$-test on the
per-map errors~\cite{neter1996applied}, as Atindama et
al.~\cite{atindama2023restoration} did. The test is paired because both
methods are applied to the same maps with the same noise, so taking the
difference map by map removes the variation between easy and hard maps. For
the same $n$ maps let $d_m = e^{(1)}_m - e^{(2)}_m$ be the difference between
the errors of the two methods on map $m$, $\bar{d}$ the mean and $S_d$ the
standard deviation of the $d_m$. The test statistic
\begin{equation}
  T = \frac{\bar{d}}{S_d / \sqrt{n}}
  \label{eq:ttest}
\end{equation}
has a $t$-distribution with $n-1$ degrees of freedom. We take $p < 0.05$ as
statistically significant.
